\begingroup
\renewcommand{\arraystretch}{1.45}
\fancyhead[L]{\small TRANSZENDENTE FUNKTIONEN}
\Needspace{10\baselineskip}
\AGMArticle{11}
Nun auch leicht die folgenden Koeffizienten der Reihe

\par\noindent\hfill\rule{0.25\linewidth}{0.4pt}\hfill\mbox{}\par
\clearpage\selectlanguage{german}\fancyhead[L]{\small ARITHMETISCH-GEOMETRISCHES MITTEL.}
\pdfbookmark[0]{FORTSETZUNG DER UNTERSUCHUNGEN ÜBER DAS ARITHMETISCH-GEOMETRISCHE MITTEL}{agm-fortsetzung}
\SourceLeaf{384}
\begin{center}\bfseries
FORTSETZUNG DER UNTERSUCHUNGEN

ÜBER DAS ARITHMETISCH-GEOMETRISCHE MITTEL.

\rule{0.08\linewidth}{0.3pt}
\end{center}
[Die weiteren Untersuchungen von \textsc{Gauss} über das arithmetisch-geometrische Mittel sind so unvollständig, nur in den Endformeln und mit so wechselnder Bezeichnung niedergeschrieben, dass fast jeder Überblick fehlen würde, wenn der Abdruck nur jene Formeln ohne begleitende Erläuterung wiedergäbe. In den folgenden Artikeln habe ich die im Nachlass gefundenen Stellen hervorgehoben und durch einen Gedankengang verbunden, der vielleicht \textsc{Gauss} dazu veranlasst hat, die betreffenden Ergebnisse unter einem erkennbaren gemeinsamen Gesichtspunkt aufzustellen.]

\Needspace{43\baselineskip}
\AGMArticle{12}
[An mehreren Stellen verwendet \textsc{Gauss} neben den in den vorhergehenden Artikeln betrachteten beiden ursprünglichen Folgen von Größen $(a,b)$, die sich auf dasselbe arithmetisch-geometrische Mittel beziehen, auch noch eine dritte Folge von Größen $(c)$. Diese hängt mit jenen durch die Gleichung $aa=bb+cc$ zusammen, die für jeden gemeinsamen Index der drei Glieder $(a,b,c)$ gilt. Nach den Festlegungen in Artikel 2 können die Gleichungen zur Berechnung der rückwärts und vorwärts folgenden, aus $a,b,c$ abgeleiteten Größen in die folgende Form gebracht werden:
\SourceLeaf{385}
\Needspace{10\baselineskip}
\[
\begin{aligned}&.\ .\quad{}^{\prime\prime}a={}^\prime a+{}^\prime c,\quad{}^\prime a=a+c,\quad a,\\&\qquad2a'=a+b,\quad2a''=a'+b',\quad.\ .\\[9pt]&.\ .\quad{}^{\prime\prime}b={}^\prime a-{}^\prime c,\quad{}^\prime b=a-c,\quad b,\\&\qquad b'b'=ab; \quad b''b''=a'b',\quad.\ .\\[9pt]&.\ .\quad{}^{\prime\prime}c{}^{\prime\prime}c=4{}^\prime a{}^\prime c,\quad{}^\prime c{}^\prime c=4ac,\quad c,\\&\qquad2c'=a-b,\quad2c''=a'-b',\quad.\ .\end{aligned}
\]
\Needspace{18\baselineskip}
Hieraus ergibt sich unmittelbar, dass, während
\Needspace{10\baselineskip}
\[
\begin{aligned}&\mathrm{M}(a,b)=\mathrm{M}(a^{\mathrm{n}},b^{\mathrm{n}})=\mathrm{M}({}^{\mathrm{n}}a,{}^{\mathrm{n}}b)\\&=\lim a^{\mathrm{m}}=\lim b^{\mathrm{m}}\quad\hbox{für}\quad m=\infty\end{aligned}
\]
gilt,
\Needspace{10\baselineskip}
\[
\begin{aligned}&\mathrm{M}(a,c)=2^n\mathrm{M}(a^{\mathrm{n}},c^{\mathrm{n}})=2^{-n}\mathrm{M}({}^{\mathrm{n}}a,{}^{\mathrm{n}}c)\\&=\lim\frac{{}^{\mathrm{m}}a}{2^m}=\lim\frac{{}^{\mathrm{m}}c}{2^m}\quad\hbox{für}\quad m=\infty\end{aligned}
\]
gilt und sich damit die dem ersten Beispiel in Artikel 3 angefügte Bemerkung, dass ${}^{\mathrm{n}}a,{}^{\mathrm{n}+1}a,{}^{\mathrm{n}+2}a$ ungefähr wie die Glieder einer geometrischen Folge mit dem Quotienten $2$ zunehmen, als allgemein gültig erweist; ferner, dass für den in Beispiel 4 von Artikel 3 berechneten Fall $a=b\surd2=c\surd2$ bei jedem $n$: ${}^{\mathrm{n}}a=2^n.a^{\mathrm{n}}$, ${}^{\mathrm{n}}b=2^n.c^{\mathrm{n}}$, ${}^{\mathrm{n}}c=2^n.b^{\mathrm{n}}$ gilt.

\Needspace{12\baselineskip}
\par Die Art, wie sich jene Größen ihren Grenzwerten nähern, ergibt sich aus den folgenden Reihenentwicklungen
\Needspace{10\baselineskip}
\[
\begin{aligned}\mathrm{M}(a,b)&=a^{\mathrm{n}}-c^{\mathrm{n}+1}-c^{\mathrm{n}+2}\\&\quad-.\;.\;.-c^{\mathrm{n}+\mathrm{m}}-.\;.\;.\end{aligned}
\]
\Needspace{10\baselineskip}
\[
\begin{aligned}\mathrm{M}(a,b)&=b^{\mathrm{n}}+c^{\mathrm{n}+1}-c^{\mathrm{n}+2}\\&\quad-.\;.\;.-c^{\mathrm{n}+\mathrm{m}}-.\;.\;.\end{aligned}
\]
\Needspace{10\baselineskip}
\[
\begin{aligned}\mathrm{M}(a,c)&=2^{-n}.{}^{\mathrm{n}}a-2^{-n-1}.{}^{\mathrm{n}+1}b-2^{-n-2}.{}^{\mathrm{n}+2}b\\&\quad-.\;.\;.-2^{-n-m}.{}^{\mathrm{n}+\mathrm{m}}b-.\;.\;.\end{aligned}
\]
\Needspace{10\baselineskip}
\[
\begin{aligned}\mathrm{M}(a,c)&=2^{-n}.{}^{\mathrm{n}}c+2^{-n-1}.{}^{\mathrm{n}+1}b-2^{-n-2}.{}^{\mathrm{n}+2}b\\&\quad-.\;.\;.-2^{-n-m}.{}^{\mathrm{n}+\mathrm{m}}b-.\;.\;.\end{aligned}
\]
\Needspace{10\baselineskip}
\[
\begin{aligned}\mathrm{M}(a,b)^2&=a^{\mathrm{n}}.a^{\mathrm{n}}-\frac12c^{\mathrm{n}}.c^{\mathrm{n}}-\frac32c^{\mathrm{n}+1}.c^{\mathrm{n}+1}\\&\quad-.\;.\;.-\frac32c^{\mathrm{n}+\mathrm{m}}.c^{\mathrm{n}+\mathrm{m}}-.\;.\;.\end{aligned}
\]
\Needspace{10\baselineskip}
\[
\begin{aligned}\mathrm{M}(a,b)^2&=b^{\mathrm{n}}.b^{\mathrm{n}}+\frac12c^{\mathrm{n}}.c^{\mathrm{n}}-\frac32c^{\mathrm{n}+1}.c^{\mathrm{n}+1}\\&\quad-.\;.\;.-\frac32c^{\mathrm{n}+\mathrm{m}}.c^{\mathrm{n}+\mathrm{m}}-.\;.\;.\end{aligned}
\]
\Needspace{10\baselineskip}
\[
\begin{aligned}\mathrm{M}(a,c)^2&=2^{-2n}.{}^{\mathrm{n}}c.{}^{\mathrm{n}}c+\frac12 2^{-2n}.{}^{\mathrm{n}}b.{}^{\mathrm{n}}b-\frac32 2^{-2n-2}.{}^{\mathrm{n}+1}b.{}^{\mathrm{n}+1}b\\&\quad-.\;.\;.\qquad-\frac32 2^{-2n-2m}.{}^{\mathrm{n}+\mathrm{m}}b.{}^{\mathrm{n}+\mathrm{m}}b-.\;.\;.\end{aligned}
\]
\Needspace{10\baselineskip}
\[
\begin{aligned}\mathrm{M}(a,c)^2&=2^{-2n}.{}^{\mathrm{n}}a.{}^{\mathrm{n}}a-\frac12 2^{-2n}.{}^{\mathrm{n}}b.{}^{\mathrm{n}}b-\frac32 2^{-2n-2}.{}^{\mathrm{n}+1}b.{}^{\mathrm{n}+1}b\\&\quad-.\;.\;.\qquad-\frac32 2^{-2n-2m}.{}^{\mathrm{n}+\mathrm{m}}b.{}^{\mathrm{n}+\mathrm{m}}b-.\;.\;.\end{aligned}
\]
\Needspace{10\baselineskip}
\[
\begin{aligned}\surd\mathrm{M}(a,b)&=\surd a^{\mathrm{n}}-\surd c^{\mathrm{n}+2}-\surd c^{\mathrm{n}+4}\\&\quad-.\;.\;.-\surd c^{\mathrm{n}+2\mathrm{m}}-.\;.\;.\end{aligned}
\]
\Needspace{10\baselineskip}
\[
\begin{aligned}\surd\mathrm{M}(a,b)&=\surd b^{\mathrm{n}}+\surd c^{\mathrm{n}+2}-\surd c^{\mathrm{n}+4}\\&\quad-.\;.\;.-\surd c^{\mathrm{n}+2\mathrm{m}}-.\;.\;.\end{aligned}
\]
\Needspace{12\baselineskip}
\SourceLeaf{386}
\Needspace{10\baselineskip}
\[
\begin{aligned}\surd\mathrm{M}(a,c)&=\surd2^{-n}.{}^{\mathrm{n}}a-\surd2^{-n-2}.{}^{\mathrm{n}+2}b-\surd2^{-n-4}.{}^{\mathrm{n}+4}b\\&\quad-.\;.\;.-\surd2^{-n-2m}.{}^{\mathrm{n}+2\mathrm{m}}b-.\;.\;.\end{aligned}
\]
\Needspace{10\baselineskip}
\[
\begin{aligned}\surd\mathrm{M}(a,c)&=\surd2^{-n}.{}^{\mathrm{n}}c+\surd2^{-n-2}.{}^{\mathrm{n}+2}b-\surd2^{-n-4}.{}^{\mathrm{n}+4}b\\&\quad-.\;.\;.-\surd2^{-n-2m}.{}^{\mathrm{n}+2\mathrm{m}}b-.\;.\;.\end{aligned}
\]
\Needspace{10\baselineskip}
\[
\begin{aligned}\displaystyle\frac\pi2.\frac{\mathrm{M}(a,b)}{\mathrm{M}(a,c)}&\displaystyle=\frac1{2^n}\log4\frac{a^{\mathrm{n}}}{c^{\mathrm{n}}}-\frac1{2^n}\log\frac{a^{\mathrm{n}}}{a^{\mathrm{n}+1}}\\&\quad-.\;.\;.-\frac1{2^{n+m}}\log\frac{a^{\mathrm{n}+\mathrm{m}}}{a^{\mathrm{n}+\mathrm{m}+1}}-.\;.\end{aligned}
\]
\Needspace{10\baselineskip}
\[
\begin{aligned}\displaystyle\frac\pi2.\frac{\mathrm{M}(a,b)}{\mathrm{M}(a,c)}&\displaystyle=\frac1{2^n}\log4\frac{b^{\mathrm{n}+1}}{c^{\mathrm{n}}}+\frac3{2^{n+1}}\log\frac{b^{\mathrm{n}+2}}{b^{\mathrm{n}+1}}\\&\quad+.\;.\;.+\frac3{2^{n+m}}\log\frac{b^{\mathrm{n}+\mathrm{m}+1}}{b^{\mathrm{n}+\mathrm{m}}}+.\;.\end{aligned}
\]
\Needspace{10\baselineskip}
\[
\begin{aligned}\displaystyle\frac\pi2.\frac{\mathrm{M}(a,c)}{\mathrm{M}(a,b)}&\displaystyle=\frac1{2^n}\log2\frac{{}^{\mathrm{n}}a}{{}^{\mathrm{n}}b}-\frac1{2^n}\log2\frac{{}^{\mathrm{n}}a}{{}^{\mathrm{n}+1}a}\\&\quad-.\;.\;.-\frac1{2^{n+m}}\log2\frac{{}^{\mathrm{n}+\mathrm{m}}a}{{}^{\mathrm{n}+\mathrm{m}+1}a}-.\;.\end{aligned}
\]
\Needspace{10\baselineskip}
\[
\begin{aligned}\displaystyle\frac\pi2.\frac{\mathrm{M}(a,c)}{\mathrm{M}(a,b)}&\displaystyle=\frac1{2^n}\log2\frac{{}^{\mathrm{n}+1}c}{{}^{\mathrm{n}}b}+\frac3{2^{n+1}}\log\frac12\frac{{}^{\mathrm{n}+2}c}{{}^{\mathrm{n}+1}c}\\&\quad+.\;.\;.+\frac3{2^{n+m}}\log\frac12\frac{{}^{\mathrm{n}+\mathrm{m}+1}c}{{}^{\mathrm{n}+\mathrm{m}}c}+.\;.\end{aligned}
\]
\Needspace{27\baselineskip}
\par In den letzten vier Gleichungen ist $\frac\pi2$ anstelle der für beständig wachsendes $m$ geltenden Grenzwerte von jeweils
\Needspace{10\baselineskip}
\[
\begin{gathered}\frac{\mathrm{M}(a^{\mathrm{m}},c^{\mathrm{m}})}{\mathrm{M}(a^{\mathrm{m}},b^{\mathrm{m}})}.\log4\frac{a^{\mathrm{m}}}{c^{\mathrm{m}}};\\[9pt]
\frac{\mathrm{M}(a^{\mathrm{m}},c^{\mathrm{m}})}{\mathrm{M}(a^{\mathrm{m}},b^{\mathrm{m}})}.\log4\frac{b^{\mathrm{m}+1}}{c^{\mathrm{m}}},\\[9pt]
\frac{\mathrm{M}({}^{\mathrm{m}}a,{}^{\mathrm{m}}b)}{\mathrm{M}({}^{\mathrm{m}}a,{}^{\mathrm{m}}c)}.\log4\frac{{}^{\mathrm{m}}a}{{}^{\mathrm{m}}b},\\[9pt]
\frac{\mathrm{M}({}^{\mathrm{m}}a,{}^{\mathrm{m}}b)}{\mathrm{M}({}^{\mathrm{m}}a,{}^{\mathrm{m}}c)}.\log2\frac{{}^{\mathrm{m}+1}c}{{}^{\mathrm{m}}b}\end{gathered}
\]
\Needspace{25\baselineskip}
gesetzt, also anstelle des Grenzwerts von $\mathrm{M}(1,\varepsilon)\log\frac4\varepsilon$ für positive Werte von $\varepsilon$, die bis auf null abnehmen. Aus jenen Gleichungen, in denen die zu logarithmierenden Werte, sobald sie alle reell sind, vom zweiten Glied der Folge an beständig bis zur Einheit hin abnehmen, wie die Beziehungen
\Needspace{10\baselineskip}
\[
\begin{gathered}\frac{a'}a=1-\frac{c'}a,\\[7pt]
\left(\frac{b'}b\right)^4=1+\frac{cc}{bb},\\[7pt]
\frac{{}^\prime a}a=2-\frac{{}^\prime b}a,\\[7pt]
\left(\frac12\frac{{}^\prime c}c\right)^4=1+\frac{bb}{cc}\end{gathered}
\]
leicht erkennen lassen, folgt zunächst, dass die gesuchte Größe völlig bestimmt ist und zum Beispiel, wenn man $a=b\surd2=c\surd2$ setzt, zwischen den Grenzen $\frac94\log2$ und $\frac{10}4\log2$ liegt. Die obigen Gleichungen gelten auch für komplexe Werte der Variablen, wenn man $n=0$ setzt und diejenigen Werte von $\log{}^{\mathrm{m}}a$, $\log{}^{\mathrm{m}}c$, $\log a$, $\log b$, $\log c$, $\log a^{\mathrm{m}}$, $\log b^{\mathrm{m}}$ zugrunde legt, die durch stetige Änderung aus den reellen Größen folgen. Nimmt man nun als ein System $(a,b,c)$ dasjenige, das der Bedingung $a=b\surd2=c\surd2$ unterliegt und aus reellen Größen besteht, und als anderes System $(\alpha,\beta,\gamma)$ dasjenige, für das $\alpha=c$, $\beta=b\surd{-1}$, $\gamma=a=b\surd2=c\surd2$ gilt und die durch Wurzelziehen zu bestimmenden Werte von $\beta^{\mathrm{m}}$ und ${}^{\mathrm{m}}\gamma$ positive Realteile erhalten, so kann man ein drittes veränderliches System $(A,B,C)$ aufstellen, das von dem einen $(a,b,c)$ zum anderen $(\alpha,\beta,\gamma)$ stetig übergeht, und zwar
\SourceLeaf{387}
so, dass bei diesem Übergang keine der Variablen ${}^{\mathrm{m}}A$, ${}^{\mathrm{m}}C$, $A$, $B$, $C$, $A^{\mathrm{m}}$, $B^{\mathrm{m}}$ den Wert null berührt oder einen negativen Realteil erhält. Die letzte jener Gleichungen gilt also, wenn $n$ gleich null gesetzt wird, sowohl für das System $(a,b,c)$ als auch für $(\alpha,\beta,\gamma)$, und da
\Needspace{10\baselineskip}
\[
\begin{gathered}\mathrm{M}(\alpha,\beta)=\tfrac12(1+\surd{-1})\mathrm{M}(a,b),\\[7pt]
\mathrm{M}(\alpha,\gamma)=\mathrm{M}(a,c)=\mathrm{M}(a,b)\end{gathered}
\]
gilt, ergibt sich $\frac{\pi}{2}\surd{-1}=\log\surd{-1}$; somit hat $\pi$ die übliche Bedeutung als Verhältniszahl der Länge des Kreisumfangs zum Durchmesser.

\Needspace{16\baselineskip}
\par Mithilfe dieser Betrachtung lässt sich auch die Richtigkeit der folgenden, von \textsc{Gauss} aufgezeichneten Bemerkung zeigen, der ich die hier verwendete Bezeichnungsweise zugrunde lege und ein von \textsc{Gauss} berechnetes Beispiel folgen lasse:]
\par \textit{Die arithmetisch-geometrischen Mittel gestalten sich anders, wenn man für eines der $b', b'', b'''\ .\ .$ den negativen Wert wählt; doch sind alle Ergebnisse in der folgenden Form enthalten:}
\Needspace{10\baselineskip}
\[
\frac{1}{\mathfrak{M}(a,b)}=\frac{1}{\mathrm{M}(a,b)}+\frac{4ik}{\mathrm{M}(a,c)}
\]
[wobei $k$ eine ganze reelle Zahl bezeichnet.]

\Needspace{38\baselineskip}
\par \textit{Beispiel für einen imaginären Wert des A. G. Mittels}
\Needspace{10\baselineskip}
\[
\begin{array}[t]{r@{\,=\,}r l r}
a&3.0000000&\log\ .\ .&0.4771213\\
b&1.0000000&&0.0000000\\[6pt]
a'&2.0000000&&0.3010300\\
b'&1.7320508&&0.2385606\\[6pt]
a''&1.8660254&&0.2709175\\
b''&1.8612098&&0.2697953\\[6pt]
a'''&1.8636176&&0.2703568\\
b'''&1.8636159&&0.2703564\\[6pt]
a''''&1.8636167&&0.2703566\\\hline
\end{array}
\]
\Needspace{10\baselineskip}
\[
\begin{array}[t]{r@{\,=\,}l l r}
a&3.0000000&\log\ .\ .&0.4771213\\
c&2.8284270&&0.4515450\\[6pt]
\frac12\mathbin{.}{}'a&2.9142\mathrm{i}35&&0.4645214\\[6pt]
\frac12\mathbin{.}{}'c&2.9129510&&0.4643332\\[6pt]
\frac14\mathbin{.}{}''a&2.9135822&&0.4644273\\[4pt]\hline
\end{array}
\]
\Needspace{43\baselineskip}
\SourceLeaf{388}
\begingroup\small\setlength{\arraycolsep}{2pt}
\Needspace{10\baselineskip}
\[
\begin{array}{r@{\,=\,}r@{}l r r r r}
a&3.0000000&&\log\ .\ .\ 0.4771213&0&&\\
b&1.0000000&&0.0000000&360^{0}&&\\[6pt]
a'&2.0000000&&0.3010300&0&&\\
b'&-1.7320508&&0.2385606&180^{0}&&\\[6pt]
a''&0.1339746&&9.1270225&0&&\\
b''&&{}+1.8612098\,i&0.2697953&90^{0}&&\\[6pt]
a'''&0.0669873&{}+0.9306049\,i&9.9698876&85^{0}&52'&58''10\\
b'''&0.3530969&{}+0.3530969\,i&9.6984089&45&0&0\\[6pt]
a''''&0.2100421&{}+0.6418509\,i&9.8295254&71^{0}&52&46.58\\
b''''&0.2836930&{}+0.6208239\,i&9.8341482&65&26&29.05\\[6pt]
a^{\mathrm{V}}&0.2468676&{}+0.6313374\,i&9.8311572&68&38&36.05\\
b^{\mathrm{V}}&0.2470649&{}+0.6324002\,i&9.8318368&68&39&37.82\\[6pt]
a^{\mathrm{VI}}&0.24699625&{}+0.6318688\,i&9.8314971&68&39&6.95\\
b^{\mathrm{VI}}&0.24699625&{}+0.6318685\,i&9.8314970&68&39&6.93\\[6pt]
a^{\mathrm{VII}}&0.24699625&{}+0.63186865\,i&9.83149705&68&39&6.94\\\hline
\end{array}
\]
\endgroup
\Needspace{10\baselineskip}
\[
\begin{aligned}&\frac{1}{\mathfrak{M}(a,b)}=+0.5365910-1.3728774\,i\\&=\frac{1}{\mathrm{M}(a,b)}+\frac{4i}{\mathrm{M}(a,c)}\end{aligned}
\]
\clearpage
\endgroup
